\section{An extremal vector for the original joint-harmonic map}

% BEGIN Y76O1
\paragraph{Y76O1. A unit extremal preimage with its image fixed.}
Let $K=U(p)\times U(q)$ act unitarily on a Hilbert space. In its
algebraic isotypic subspace for one irreducible lowest $K$-type
let $v$ be a unit smooth vector. Suppose a linear map $A$ on the
smooth part intertwines the $K$-action with an action on its image.
Use the original torus ordering, in which the order of the negative
block is reversed. Write $\lambda_*$ for the resulting extremal
weight. Number the rows from zero and define the height of a weight by
\[
       H(\nu)=-\sum_{a<p}a\nu_a+\sum_{p\le a<p+q}a\nu_a .
\]
Suppose $Av$ has weight $\mu$ and every nonzero weight vector
in the span of its $K$-translates has height at most $H(\mu)$.
Then there is a smooth unit vector $v'$ of weight $\lambda_*$
with $Av'=Av$.

Take the cyclic subspace $U$ generated by $v$. It is finite
dimensional: an element of the algebraic isotypic sum is contained
in a sum of finitely many copies of the irreducible $K$-type.
Smooth vectors remain smooth under $K$, by conjugating the orbit
map, and hence under finite linear combinations. Split $U$ orthogonally
as $J\oplus J^\perp$ inside $U$, where $J=\ker(A|_U)$.
Both subspaces are $K$-stable. Write $v=v_J+v_1$ accordingly.
If $v_1\ne0$, equivariance and injectivity on $J^\perp$ force
$v_1$ to have weight $\mu$. Every nonzero invariant summand
contains a vector of weight $\lambda_*$: a finite sum of copies
of the given type is completely reducible by invariant orthogonal
complements, and a permutation matrix reverses the negative
block of a highest vector. Apply this to $J^\perp$. Its extremal
vector has nonzero image, whence $H(\lambda_*)\le H(\mu)$.
Conversely every weight of this type is below its highest weight
for this ordering. Its difference from $\lambda_*$, subtracted
in the order $\lambda_*-\mu$, is a sum of nonnegative multiples
of roots $e_a-e_b$ with $a<b$ in the positive block or $a>b$
in the negative block. Each root has strictly positive height.
Thus $\mu=\lambda_*$. If $v_J=0$ take $v'=v_1$.
If $v_J\ne0$, take a unit extremal vector $e\in J$ and put
$v'=v_1+\|v_J\|e$. Orthogonality gives $\|v'\|=1$, and
$Av'=Av$. This also covers $v_1=0$ and a wholly zero map.
In particular the construction requires no positive kernel
dimension.
% END Y76O1

% BEGIN Y76O2
\paragraph{Y76O2. The product of minors is extremal in its cyclic image.}
At an infinite place $\tau$ of the original base theta pair,
write $K=U(p)\times U(q)$ in its given orthonormal coordinates
and $\Delta$ for the distinguished polynomial. In the Fock
polynomial action at $\tau$ it is a torus eigenvector of maximal
$H$ among the weights in its $K$-span.

Here the vacuum action is a scalar character $\chi(k)$ times
the substitution on the creation variables. The complex structure
and the Gaussian vacuum line are preserved by $K$. An element
supported at $\tau$ acts in that archimedean oscillator factor.
Conjugation of a creation operator by this action is the
corresponding creation operator in the transformed direction,
by Heisenberg covariance. Applying it once per letter gives
the substitution formula for every monomial. The substitution
on the two row blocks is by $k_+$ and $\bar k_-$ in the convention
that the image of a letter uses its column of the substituted
matrix. The scalar on the vacuum is a continuous unitary
character; its restriction to either special unitary group
is trivial since that group is generated by commutators.
It is therefore a product of determinant powers on the
nonempty blocks. On the diagonal its weight is constant on each
block, say $c_a$. An empty block contributes zero.

Each factor of $\Delta$ at $\tau$ is a minor using the first
rows in its sign block; its columns are those specified in
the native polynomial. Let $e_a$ be the row degrees of their
product, with repetitions given by its exponents. Every
determinant monomial uses each selected row once. Thus the
diagonal weight is $\mu_a=c_a+s_ae_a$, where $s_a=1$ in the
positive block and $s_a=-1$ in the negative block. Under
substitution by a block matrix, Cauchy--Binet expands a minor
of size $b$ into minors with $b$ distinct row indices in the
same block. The sum of these indices is at least the sum of
the first $b$ indices there. Hence every monomial of every
translate, and therefore of every linear combination, has
weighted row degree $\sum_a a d_a\ge\sum_a a e_a$. Other
places are unchanged. If such a combination is a nonzero
torus eigenvector of weight $\nu$, read a single nonzero
monomial. Comparing its characters on all independent angles
gives $\nu_a=c_a+s_a d_a$. Consequently
$H(\nu)=H(c)-\sum_a a d_a\le H(\mu)$. Constant minors and
zero exponents satisfy the same calculation.
% END Y76O2

% BEGIN Y76O3
\paragraph{Y76O3. The oscillator extension from the weak characterization.}
For $R=R_3.\mathrm{toRealDat}$, the original weak characterization
at $F_3$ and $R.\mathrm{ExAt}$ imply $\mathrm{IoExt}(R)$.
The needed inputs at the second stage are its level-character,
model and point-theta laws. The first two come immediately
from the weak base. For the third, apply each conditional weak
oscillator and theta equation to the corresponding $\mathrm{ExAt}$
witness, as in the first paragraph of T76R1. The proof makes
no assumption about the dimension of a space of automorphic lifts.

Fix $i$ and the base tuple. Applying the polynomial oscillator
map to $1$ gives the original finite theta vector with its real
vacuum factors. Its nonvanishing and injectivity on polynomials
use just the first part of P74I2: the level frames are self-dual,
the intrinsic rank inequality $\operatorname{rank} W_i\le
\operatorname{rank} V_i$ supplies the full residual block, and
one may take identity blocks at the marked residues and the
specified opposite conditions at the other half. The resulting
finite test is nonzero. Its Schr\"odinger realization on the
unit vacuum is injective by the distinct translation characters
on a self-dual lattice. On the real factors the Fock
annihilators kill this input; commuting them past the creation
variables differentiates a polynomial. For a nonzero polynomial
differentiate through a maximal-degree monomial to obtain a
nonzero scalar times the nonzero input. This proves injectivity
for the precise creation map, including its given orthonormal
source coordinates.

The point-oscillator law at the base tuple gives unit vectors
$v_\rho$ of the smooth lowest types and a unit coefficient vector
$w$ with $\mathrm{ioP}_i(\bigotimes_\rho v_\rho,w)=\Delta$.
Fix all these variables except $v_\tau$. Their contraction
defines a linear map $A_\tau$ on the smooth lowest type at
$\tau$. The compact covariance in the same law, followed by
the polynomial action of Y76O2, shows that it intertwines:
both expressions agree after the creation map and that map
is injective. Y76O1 and Y76O2 replace $v_\tau$ by an
extremal unit vector with precisely the same contracted image.
If the image vanishes the kernel case in Y76O1 applies as well.
Proceed through the finite set of infinite places, keeping
$w$ unchanged. Previously treated factors remain unit and
extremal. The resulting family is the original
$\mathrm{IsExtVec}$ and pairs to the exact $\Delta$ with
that unit $w$, proving $\mathrm{IoExt}$ at every $i$.

Consequently in X76R2 and X76R3 the oscillator extension follows
from the already specified weak characterization and $\mathrm{ExAt}$.
On the selected column with the cover field of U76W3, those
two hypotheses, the geometric datum, rank bound and
$\mathrm{LUnif}$ give both row maps. With the other fields
of $\mathrm{TiesB}$ the original rational tests and lattices
then follow from X76R3 and T76R4.
% END Y76O3
